% Materiał z AF 25; fragment do rozdziału o przestrzeniach Montela.
\subsection{Konsekwencje dla topologii przestrzeni sprzężonych}

\begin{uwaga}[Zupełność i lokalna wypukłość]
\label{rem:tabela-struktur25}
$\mathcal E(\Omega)$ i $\mathcal S(\R^d)$ są nuklearnymi
przestrzeniami Frécheta, a $\mathcal D(\Omega)$ jest nuklearną,
zupełną, niemetryzowalną przestrzenią LF
(Twierdzenia~\ref{thm:D-complete18}, \ref{thm:D-nonmetric18}).
Wszystkie trzy są przestrzeniami Montela
(Twierdzenie~\ref{thm:montel-test23}).
Nuklearność $\mathcal E$ i $\mathcal S$ wynika odpowiednio
z Twierdzenia~\ref{thm:nuclear-cinfty23}
i Wniosku~\ref{cor:nuclear-schwartz23}.
Trwałość nuklearności względem podprzestrzeni, ilorazów
i przeliczalnych granic induktywnych udowodniono w
Twierdzeniach~\ref{thm:nuclear-subquot23} i~\ref{thm:nuclear-indlim23};
nuklearność $\mathcal D_K$ i $\mathcal D$ wynika z
Twierdzenia~\ref{thm:nuclear-D23}. Nie należy mylić nuklearności
z własnością Montela: te dwie własności były w skrypcie
uzasadniane oddzielnie.
\end{uwaga}
\begin{twierdzenie}[Zbieżność ciągów w przestrzeniach dualnych]
\label{thm:weak-strong-sequences25}
Niech $X$ będzie beczkową przestrzenią lokalnie wypukłą
Hausdorffa nad $\K$, w której każdy zbiór ograniczony
ma zwarte domknięcie (w szczególności przestrzenią Montela).
Dla ciągu $(T_n)$ w $X'$ i $T\in X'$ następujące
warunki są równoważne: $T_n\to T$ w topologii
$\sigma(X',X)$ oraz $T_n\to T$ w topologii $\beta(X',X)$.
Stwierdzenie stosuje się w szczególności do
$X=\mathcal D(\Omega)$, $X=\mathcal E(\Omega)$ i
$X=\mathcal S(\R^d)$.
\end{twierdzenie}
\begin{proof}
Założenia oznaczają, że $X$ jest przestrzenią Montela.
Zbieżność słaba-* pociąga za sobą silną na mocy wniosku
\ref{cor:montel-dual-seq23}. Odwrotna implikacja wynika
z ograniczoności zbiorów jednoelementowych.
Zastosowania do $\mathcal D$, $\mathcal E$ i $\mathcal S$
wynikają z twierdzenia \ref{thm:montel-test23}.
\end{proof}


\begin{uwaga}[Zbieżność ciągów a równość topologii]
\label{rem:sequence-net25}
Twierdzenie~\ref{thm:weak-strong-sequences25} nie oznacza
równości topologii $\sigma$ i $\beta$.
W dowodzie używamy punktowej ograniczoności całej
rodziny wyrazów ciągu zbieżnego. Dla sieci zbieżnej
punktowo nie wynika stąd ograniczoność wszystkich
jej wyrazów.

Dla każdej z trzech rozważanych przestrzeni $X$
istnieje ograniczony zbiór $B_0$ zawierający nieskończenie
wiele liniowo niezależnych funkcji. Dla $\mathcal D$
jest to zbiór skonstruowany w dowodzie Twierdzenia~
\ref{thm:weak-strong-different19}; dla $\mathcal E$
i $\mathcal S$ korzystamy z jego obrazu przez ciągłe
włożenie z Twierdzenia~\ref{thm:zanurzenia25}.
Niech $I$ będzie zbiorem par $(F,n)$, gdzie
$F\subset X$ jest skończony, $n\in\N$, uporządkowanym
przez inkluzję pierwszego czynnika i zwykły porządek
drugiego. Dla każdej pary wybieramy
$u\in B_0\setminus\operatorname{lin}F$.
Skończenie wymiarowa podprzestrzeń $\operatorname{lin}F$
jest domknięta, więc twierdzenie Hahna--Banacha
zapewnia $T_{F,n}\in X'$ z
$T_{F,n}|_{\operatorname{lin}F}=0$ i $T_{F,n}(u)=n$.
Dla ustalonego $x\in X$ od indeksu $(\{x\},1)$
wszystkie wartości $T_{F,n}(x)$ są zerowe.
Sieć zbiega więc słabo-* do zera. Jednak
$\sup_{v\in B_0}|T_{F,n}(v)|\ge n$, więc
nie zbiega silnie do zera. Topologia silna
jest zatem ściśle silniejsza w każdym z trzech przypadków.
\end{uwaga}
